A further change of coordinate lowers the degrees of the ordinary
construction. Replace the supplied \(y\) by a supplied \(e\), and
compute \(y=V+e\). The auxiliary factor becomes
\begin{equation}
 N_A^{\mathrm{gap}}
  =U\bigl(V^2-(V+e)^2\bigr)+(V+e)^2
  =V^2-(U-1)e(2V+e).
 \label{eq:gap-auxiliary}
\end{equation}
Keep the other factors and residuals unchanged. Thus the new product is
\[
 P_{70}^{\mathrm{gap}}
   =N_0N_1N_A^{\mathrm{gap}}N_kN_tN_s-1=0.
\]
and the new six-equation system is
\begin{equation}\boxed{\begin{aligned}
 N_0=N_1=N_A^{\mathrm{gap}}=N_k&=1,\\
 (K+w)C&=F+z(q-1),\\
 I^2&=U.
\end{aligned}}\label{eq:gap-system61}\end{equation}
It replaces only the third equation of
\eqref{eq:system60} by \(N_A^{\mathrm{gap}}=1\).
The sum of its six squared residuals gives a third presentation.
The circuit needs just one additional addition: its existing
\(y\)-square now squares the computed value \(V+e\).

The signed-domain equivalence is immediate from the inverse integer
substitutions \(e=y-V\) and \(y=V+e\). The nine coordinates
\(J,F,\alpha,z,i,s,w,\eta,\zeta\) remain positive, while
\(h,f,T,\tau,e,\gamma\) may be arbitrary integers. Soundness therefore
follows from Section~\ref{sec:arithmetic-proof}. Positive-domain
completeness needs one extra check. The canonical positive auxiliary
Pell point used by the blank compiler has \(V,y>0\), \(U>1\),
\(c>2R\), and \(V\equiv-R\pmod c\). Hence
\(V\ge c-R>1\), and its norm equation gives
\[
 (U-1)(y^2-V^2)=V^2-1>0.
\]
Thus \(y>V\) and the canonical gap \(e=y-V\) is strictly positive.
This proves the same existence equivalence in both stated domains.

\begin{center}
\begin{tabular}{lrrr}
\toprule
Presentation & Unknowns & Operations & Exact degree\\
\midrule
Six equations & \(15\) positive & 61 & maximum 22\\
Product & \(15\) positive & 70 & 93\\
Sum of squares & \(15\) positive & 78 & 44\\
Six equations & \(33\) integer & 124 & maximum 44\\
Product & \(33\) integer & 133 & 179\\
Sum of squares & \(33\) integer & 141 & 88\\
\bottomrule
\end{tabular}
\end{center}

For exact degree certificates, set \(Q_0=\beta J\),
\(k_0=\eta+\zeta\), and
\(T_2=w(Q_0-F-\alpha)-z_0Q_0\). The auxiliary leader is
\(-2ek_0w^2s^3Q_0^{11}Tf^3\); the six factor degrees are
\(22,18,22,7,2,22\). Consequently
\begin{equation}
 (P_{70}^{\mathrm{gap}})_{93}
  =-16Q_0^{52}h\gamma i^2k_0^8w^7s^{15}Tf^3eT_2.
 \label{eq:gap-product-leader}
\end{equation}
The coefficient of
\(J^{53}h\gamma i^2\eta^8w^7s^{15}Tf^3ez_0\)
is \(16\beta^{53}\ne0\). The degree-22 leaders of the first,
auxiliary and strong residuals are
\[
 -w^2k_0^2s^4Q_0^{14},\qquad
 -2ek_0w^2s^3Q_0^{11}Tf^3,\qquad
 i^2k_0^4s^4Q_0^{12}.
\]
Their squared sum is nonzero, proving degree 44 for the sum of squares.

For the integer conversion, give weight two to the nine positive
coordinates and weight one to the other six. The residual weights are
\(44,35,39,13,4,44\). Nine independent substitutions
\eqref{eq:pos} add 63 gates and give 33 integer unknowns. The product
has degree 179, while the first and strong squared residuals give
degree 88 for the sum of squares. The disjoint homogeneous quadratic
substitutions preserve the nonzero leading forms. Complete circuits,
ledgers and checks are emitted by \code{arithmetic/build_gap.py}.
